% examples/report.rmd の図「reportmd の処理の流れ」の元データ。
% 作り直すときは、このディレクトリで次を実行する：
%   xelatex pipeline.tex && pdftoppm -r 300 -png -singlefile pipeline.pdf pipeline
\documentclass[border=6pt]{standalone}
\usepackage{xeCJK}
\setCJKmainfont{IPAexGothic}
\setmainfont{IPAexGothic}
\usepackage{tikz}
\usetikzlibrary{positioning}
\begin{document}
\begin{tikzpicture}[
  node distance=7mm,
  stage/.style={draw, thick, rounded corners=2pt, align=center,
               minimum height=12mm, minimum width=20mm, fill=black!5},
  file/.style={draw, thick, align=center, minimum height=12mm, minimum width=14mm},
  arrow/.style={->, thick}
]
  \node[file] (rmd) {.rmd};
  \node[stage, right=of rmd] (pre) {前処理\\\footnotesize preprocessor.rb};
  \node[stage, right=of pre] (pandoc) {Pandoc};
  \node[stage, right=of pandoc] (filter) {フィルタ\\\footnotesize filter.lua};
  \node[stage, right=of filter] (xelatex) {XeLaTeX};
  \node[file, right=of xelatex] (pdf) {PDF};
  \node[file, above=of pandoc] (setting) {setting.json};

  \draw[arrow] (rmd) -- (pre);
  \draw[arrow] (pre) -- (pandoc);
  \draw[arrow] (pandoc) -- (filter);
  \draw[arrow] (filter) -- (xelatex);
  \draw[arrow] (xelatex) -- (pdf);
  \draw[arrow, dashed] (setting) -- (pandoc);
  \draw[arrow, dashed] (setting) -| (xelatex);
\end{tikzpicture}
\end{document}
